\section{Equity Impact: In-Between Communities}
\label{sec:equity}

\subsection{The In-Between Problem}

The 12 T1 corridors analyzed in this paper pass through hundreds of communities that are not major metropolitan areas and currently have no intercity transit option. These are the communities that the US bus deregulation of 1982 \citep{BusRegulatoryReform1982} theoretically left to the market, and that the market has systematically abandoned \citep{Schwieterman2019}. They are called \emph{in-between} communities in this paper: places between the metro pairs that anchor T1 corridors, close enough to a T1 route to access a hub stop but too small to anchor a direct O-D market.

In-between communities are characterized by three features relevant to T1 bus equity:
\begin{enumerate}
  \item \textbf{High zero-vehicle household rates.} American Community Survey Table B08201 \citep{ACS_B08201_2022} shows that communities along T1 corridors outside major metros have zero-vehicle household rates of 15--28\%, compared to a national average of 8\%. Agricultural counties along I-35 in Texas and Oklahoma average 19\% zero-vehicle households; Mississippi counties along I-20 average 24\%.

  \item \textbf{High C3 scores.} The ROUTE scoring dimension C3 (Economic Opportunity Access deficit) measures the gap between a community's economic productivity and the national mean, weighted by transit dependency. In-between T1 communities score C3 = 6.5--8.5 on the 10-point ROUTE rubric \citep{ROUTE_A1}, significantly above the corpus mean of 4.8. These communities have the most to gain from access to regional employment, healthcare, and educational centers.

  \item \textbf{No current intercity transit.} Greyhound's post-2020 network covers fewer than 300 intermediate stops nationally; T1/T2 hub locations (approximately 50 per the F.1 network \citep{ROUTE_F1}) would bring intercity transit access to communities that have not had scheduled bus service for 10--20 years.
\end{enumerate}

\subsection{Representative In-Between Communities}

\paragraph{Coalinga, CA (I-5, midway LA--SF, population 17,000).}

Coalinga is an agricultural service center in the San Joaquin Valley, 75 miles from Fresno and 115 miles from Bakersfield. The nearest current Greyhound stop is Fresno (52 miles by highway). A T1/T2 hub at the I-5/CA-33 interchange --- already a major truck rest area with fueling infrastructure --- would place intercity transit access within 15 miles of approximately 35,000 residents in the surrounding Fresno County agricultural region.

Zero-vehicle household rate: 22\% (ACS 2022, Fresno County agricultural communities \citep{ACS_B08201_2022}). These are predominantly agricultural workers and their families. Without a vehicle, accessing Fresno for medical appointments, Bakersfield for court appearances, or San Francisco for family requires either a 2-hour drive with a neighbor or no trip at all.

\paragraph{Effingham, IL (I-70/I-57, population 12,000).}

Effingham sits at the intersection of I-70 (Kansas City--Indianapolis) and I-57 (Chicago--Memphis), making it a natural T1/T1 diamond hub candidate. It is equidistant between St.\ Louis (105 miles) and Indianapolis (130 miles), and 210 miles south of Chicago. Greyhound withdrew its Effingham stop in 2020 \citep{Schwieterman2019}. A T1 hub at the I-70/I-57 diamond would serve approximately 45,000 residents within a 30-mile radius, including communities in Wayne, Clay, Cumberland, and Jasper counties.

Zero-vehicle household rate: 16\% (ACS 2022, Effingham County and adjacent counties \citep{ACS_B08201_2022}). The Effingham region has disproportionately elderly population (22\% over 65, vs.\ 17\% national) and disproportionately low income (median household income \$48,200, vs.\ \$74,600 national). Healthcare access is the primary unmet need: the nearest Level II trauma center is 105 miles away in St.\ Louis.

\paragraph{Durant, OK (I-35, population 18,000).}

Durant is a county seat in Bryan County, 90 miles north of Dallas and 150 miles south of Oklahoma City on I-35. It hosts Southeastern Oklahoma State University (4,000 enrollment). Greyhound withdrew its Durant stop in 2020. The community has no intercity transit option. A T1/T2 hub at the existing I-35/US-70 interchange rest area would serve approximately 55,000 residents in Bryan County and adjacent Atoka and Marshall counties.

Zero-vehicle household rate: 21\% (ACS 2022, Bryan County \citep{ACS_B08201_2022}). University students without vehicles, agricultural workers, and elderly residents are the primary zero-vehicle populations. Dallas employment access is the primary unmet demand: the 90-mile I-35 trip at T1 bus speed is 1.5 hours, a commutable distance for weekly or bi-weekly work travel.

\subsection{Equity Metrics: Cost and Access}

\paragraph{Fare cost.} T1 express bus targeted at \$0.12/mile produces the following representative fares:
\begin{itemize}
  \item Coalinga $\rightarrow$ San Francisco (175 miles): \$21.00 one-way
  \item Effingham $\rightarrow$ Chicago (210 miles): \$25.20 one-way
  \item Durant $\rightarrow$ Dallas (90 miles): \$10.80 one-way
\end{itemize}

These fares are comparable to Greyhound economy pricing (\$0.09--0.14/mile) and dramatically below air (\$0.30--0.60/mile at short-haul distances after fees). For a zero-vehicle household accessing medical care in a regional center, \$21.00 round-trip is a feasible expenditure; rental car plus gas for the same trip costs \$80--120.

\paragraph{Healthcare access.} The Medicaid Non-Emergency Medical Transportation (NEMT) program provides transportation for Medicaid beneficiaries to medical appointments at an average cost of \$60/trip \citep{BTS_intermodal2023}. T1 bus service to in-between communities would convert NEMT demand (dedicated van, high marginal cost) into scheduled transit demand (\$21 round-trip, marginal public cost near zero). For communities like Effingham with 15--20\% Medicaid enrollment, this represents a significant public health budget substitution.

\paragraph{C3 and equity multiplier.} The F.1 analysis \citep{ROUTE_F1} showed that T1/T1 hub locations align disproportionately with high-C3 communities. F.2 confirms this pattern at the in-between stop level: the 50 T1/T2 regional stops in the F.1 network have a mean C3 score of 7.1, versus 4.8 for the full corridor corpus. In-between communities are not randomly distributed along T1 corridors; they are concentrated at the points of maximum economic-opportunity deficit, precisely where transit access has the highest marginal value.

\subsection{Comparison to Existing Programs}

The Essential Air Service program spends approximately \$175/passenger \citep{EAS_DOT2023} to maintain commercial air service to approximately 175 communities with populations in the 2,000--30,000 range. Many EAS communities are in exactly the income and vehicle-ownership categories that characterize in-between T1 communities. The proposed Essential Intercity Transportation designation, at a subsidy cap of \$15/passenger, delivers the same connectivity policy at 8.6 cents per EAS dollar --- a 91.4\% cost reduction in the public subsidy per passenger served, with no requirement for airport infrastructure. Applying EIT to the 20--25 most isolated in-between T1 communities would cost \$30--50 million annually, a rounding error against the EAS appropriation of approximately \$360 million/year \citep{EAS_DOT2023}.
